\subsection{Nakayama's lemma}

We now state a very useful criterion for determining when a module over a
\emph{local} ring is zero.


\begin{lemma}[Nakayama's lemma] \label{nakayama} If $R$ is a local ring with
maximal ideal
$\mathfrak{m}$. Let $M$ be a finitely generated $R$-module.  If
$\mathfrak{m}M = M$, then $M = 0$.
\end{lemma}

Note that $\mathfrak{m}M$ is the submodule generated by products of
elements of $\mathfrak{m}$ and $M$.

\begin{remark}
Once one has the theory of the tensor product, this equivalently states that
if $M$ is finitely generated, then
\[ M \otimes_R R/\mathfrak{m} = M/\mathfrak{m}M \neq 0.  \]
So to prove that a finitely generated module over a local ring is zero, you
can reduce to studying the reduction to $R/\mathfrak{m}$. This is thus a very
useful criterion.
\end{remark}

Nakayama's lemma highlights why it is so useful to work over a local ring.
Thus, it is useful to reduce questions about general rings to questions about
local rings.
Before proving it, we note a corollary.

\begin{corollary}
Let $R$ be a local ring with maximal ideal $\mathfrak{m}$, and $M$ a finitely
generated module. If $N \subset M$ is a submodule such that $N +
\mathfrak{m}N =
M$, then $N=M$.
\end{corollary}
\begin{proof}
Apply Nakayama above (\cref{nakayama}) to $M/N$.
\end{proof}


We shall prove more generally:

\begin{proposition} 
Suppose $M$ is a finitely generated $R$-module, $J \subset R$ an ideal.
Suppose $JM = M$. Then there is $a \in 1+J$ such that $aM = 0$.
\end{proposition} 

If $J$ is the maximal ideal of a local ring, then $a$ is a unit, so that $M=0$.

\begin{proof}
Suppose $M$ is generated by $\left\{x_1, \dots, x_n\right\} \subset M$. This
means that every element of $M$ is  a linear combination of elements of
$x_i$. However, each $x_i \in JM$ by assumption. In particular, each
$x_i$ can be written as
\[ x_i = \sum a_{ij} x_j, \ \mathrm{where} \ a_{ij} \in \mathfrak{m}.  \]
If we let $A$ be the matrix $\left\{a_{ij}\right\}$, then $A$ sends the
vector  $(x_i)$ into itself. In particular, $I-A$ kills
the vector $(x_i)$.

Now $I-A$ is an $n$-by-$n$ matrix in the ring $R$. We could, of course,
reduce everything modulo $J$ to get the identity; this is
because $A$ consists of elements of $J$. It follows that the
determinant must be congruent to $1$ modulo $J$.

In particular, $a=\det (I - A)$ lies in $1+J$.  
Now by familiar linear algebra, $aI$ can be represented as the product of $A$
and the matrix of cofactors; in particular, $aI$ annihilates the vector
$(x_i)$, so that $aM=0$.
\end{proof}

Before returning to the special case of local rings, we observe the following
useful fact from ideal theory:

\begin{proposition}  \label{idempotentideal}
Let $R$ be a commutative ring, $I \subset R$ a finitely generated ideal such that $I^2 = I$.
Then $I$ is generated by an idempotent element.
\end{proposition} 
\begin{proof} 
We know that there is $x \in 1+I$ such that $xI =0$. If $x = 1+y, y \in I$, it
follows that 
\[ yt = t  \]
for all $t \in I$. In particular, $y$ is idempotent and $(y) = I$.
\end{proof} 

\begin{exercise} 
\rref{idempotentideal} fails if the ideal is not finitely generated.
\end{exercise} 

\begin{exercise} 
Let $M$ be a finitely generated module over a ring $R$. Suppose $f: M \to M$
is a surjection. Then $f$ is an isomorphism. To see this, consider $M$ as a
module over $R[t]$ with $t$ acting by $f$; since $(t)M = M$, argue that there
is a polynomial $Q(t) \in R[t]$ such that $Q(t)t$ acts as the identity on
$M$, i.e. $Q(f)f=1_M$.
\end{exercise} 

\begin{exercise}
Give a counterexample to the conclusion of Nakayama's lemma when the module is
not finitely generated.
\end{exercise}
\begin{exercise}
Let $M$ be a finitely generated module over the ring $R$. Let $\mathfrak{I}$
be the Jacobson
radical of $R$ (cf. \rref{Jacobson}). If $\mathfrak{I} M = M$,
then $M =
0$.
\end{exercise}

\begin{exercise}[A converse to Nakayama's lemma]
Suppose conversely that $R$ is a ring, and $\mathfrak{a} \subset R$ an ideal
such that $\mathfrak{a} M \neq M$ for every nonzero finitely generated
$R$-module. Then $\mathfrak{a}$ is contained in every maximal ideal of $R$.
\end{exercise}


\begin{exercise}
Here is an alternative proof of Nakayama's lemma. Let $R$ be local with
maximal ideal $\mathfrak{m}$, and let $M$ be a finitely generated module with
$\mathfrak{m}M = M$. Let $n$ be the minimal number of generators for $M$. If
$n>0$, pick generators $x_1, \dots, x_n$. Then write $x_1 = a_1 x_1 + \dots +
a_n x_n$ where each $a_i \in \mathfrak{m}$. Deduce that $x_1$ is in the
submodule generated by the $x_i, i \geq 2$, so that $n$ was not actually
minimal, contradiction.
\end{exercise}

Let $M, M'$ be finitely generated modules over a local ring $(R,
\mathfrak{m})$, and let $\phi: M \to M'$ be a homomorphism of modules. Then
Nakayama's lemma gives a criterion for $\phi$ to be a surjection: namely, the
map $\overline{\phi}: M/\mathfrak{m}M \to M'/\mathfrak{m}M'$ must be a surjection.
For injections, this is false. For instance, if $\phi$ is  multiplication by any element of
$\mathfrak{m}$, then $\overline{\phi}$ is zero but $\phi$ may yet be injective.
Nonetheless, we give a criterion for a map of \emph{free} modules over a local ring to
be a \emph{split}  injection.

\begin{proposition} \label{splitcriterion1}
Let $R$ be a local ring with maximal ideal $\mathfrak{m}$. Let $F, F'$ be two
finitely generated free $R$-modules, and let $\phi: F \to F'$ be a homomorphism. 
Then $\phi$ is a split injection if and only if the reduction $\overline{\phi}$
\[ F/\mathfrak{m}F \stackrel{\overline{\phi}}{\to} F'/\mathfrak{m}F'  \]
is an injection. 
\end{proposition}
\begin{proof} 
One direction is easy. If $\phi$ is a split injection, then it has a left
inverse
$\psi: F' \to F$ such that $\psi \circ \phi = 1_F$. The reduction of $\psi$ as a
map $F'/\mathfrak{m}F' \to F/\mathfrak{m}F$ is a left inverse to
$\overline{\phi}$, which is thus injective.

Conversely, suppose $\overline{\phi}$ injective. Let $e_1, \dots, e_r$ be a
``basis'' for $F$, and let $f_1, \dots, f_r$ be the images under $\phi$ in
$F'$. Then the reductions $\overline{f_1}, \dots, \overline{f_r}$ are linearly
independent in the $R/\mathfrak{m}$-vector space $F'/\mathfrak{m}F'$. Let us
complete this to a basis of $F'/\mathfrak{m}F'$ by adding elements
$\overline{g_1}, \dots, \overline{g_s} \in F'/\mathfrak{m}F'$, which we can
lift to elements $g_1, \dots, g_s \in F'$. It is clear that $F'$ has rank $r+s $
since its reduction $F'/\mathfrak{m}F'$ does.

We claim that the set $\left\{f_1, \dots, f_r, g_1, \dots, g_s\right\}$ is a
basis for $F'$. Indeed, we have a map 
\[ R^{r+s} \to F'  \]
of free modules of rank $r+s$. It can be expressed as an $r+s$-by-$r+s$ matrix
$M$; we need to show that $M$ is invertible. But if we reduce modulo
$\mathfrak{m}$, it is invertible since the reductions of $f_1, \dots, f_r,
g_1, \dots, g_s$ form a basis of $F'/\mathfrak{m}F'$. 
Thus the determinant of $M$ is not in $\mathfrak{m}$, so by locality it is
invertible. 
The claim about $F'$ is thus proved.

We can now define the left inverse $F' \to F$ of $\phi$. Indeed, given $x \in F'$,
we can write it uniquely as a linear combination $\sum a_i f_i + \sum b_j g_j$
by the above. We define $\psi(\sum a_i f_i + \sum b_j g_j) = \sum a_i e_i \in
F$. It is clear that this is a left inverse
\end{proof} 

We next note  a slight strenghtening of the above result, which is sometimes
useful. Namely, the first module does not have to be free.
\begin{proposition} 
Let $R$ be a local ring with maximal ideal $\mathfrak{m}$. Let $M, F$ be two
finitely generated $R$-modules with $F$ free, and let $\phi: M \to F'$ be a homomorphism. 
Then $\phi$ is a split injection if and only if the reduction $\overline{\phi}$
\[ M/\mathfrak{m}M \stackrel{\overline{\phi}}{\to} F/\mathfrak{m}F  \]
is an injection. 
\end{proposition}
It will in fact follow that $M$ is itself free, because $M$ is projective (see
\cref{} below) as it is a direct summand of a free module.
\begin{proof} 
Let $L$ be a ``free approximation'' to $M$.
That is, choose a basis $\overline{x_1}, \dots, \overline{x_n}$ for $M/\mathfrak{m}M$ (as an $R/\mathfrak{m}$-vector
space) and lift this to elements $x_1, \dots, x_n \in M$. Define a map
\[ L = R^n \to M  \]
by sending the $i$th basis vector to $x_i$.
Then $L/\mathfrak{m} L \to M/\mathfrak{m}M$ is an isomorphism.
By Nakayama's lemma,
$L \to M$ is surjective.

Then the composite map
$L \to M \to F$ is such that the $L/\mathfrak{m}L \to F/\mathfrak{m}F$ is injective, so
$L \to F$ is a split injection (by \cref{splitcriterion1}).
It follows that we can find a splitting $F \to L$, which when composed with $L
\to M$ is a splitting of $M \to F$.
\end{proof} 

\begin{exercise} 
Let $A$ be a local ring, and $B$ a ring which is finitely generated and free as an
$A$-module. Suppose $A \to B$ is an injection. Then $A \to B$ is a \emph{split
injection.} (Note that any nonzero morphism mapping out of a field is
injective.) 
\end{exercise} 

